\section{Strict interlacing for split graphs}\label{sec:interlacing}

Write $E_{k,n}$ in this section for the Ehrhart polynomial of the graph
$K_{1^k,n}$, with $k$ singleton parts and one part of size $n$.
K\"olbl's interlacing conjecture asks what happens when one adds a
singleton, keeping $n$ fixed. His Theorem~3.3(a,d) gives the first two
steps~\cite{kolbl}. The following theorem has no upper bound on $n$.

\begin{theorem}\label{thm:split}
For $1\le k\le20$ and every positive integer $n$, the zeros of
$E_{k,n}$ and $E_{k+1,n}$ are simple, have real part $-1/2$, and strictly
interlace when ordered by imaginary part. In particular the simple-root
canonical-line property holds through $21$ singleton parts.
\end{theorem}

\begin{proof}
Let $r=k-1$, $d=n+r$. The formula of Theorem~\ref{thm:gamma}
specializes to
\[
 g_j=\binom{2j}{j}\sum_{\ell=0}^j
           \binom n\ell\binom r{2j-\ell}\qquad (0\le j\le r).
\]
Only the $n$-part can contain more than half of a set of size $2j$,
for $j\ge1$. We set $g_j=0$ for $2j>d$, as the formula already does.
Let $C_m$ be the cross-polytope Ehrhart polynomial and put
\[
 Q_m(y)=m!\,i^{-m}C_m((iy-1)/2).
\]
Then $Q_0=1$, $Q_1=y$ and
$Q_{m+1}=yQ_m-m^2Q_{m-1}$. This recurrence also follows by
differentiating the generating function
$\sum_{m\ge0}C_m(x)u^m=(1+u)^x/(1-u)^{x+1}$.
The gamma expansion, the Ehrhart series
$(1+t)^m/(1-t)^{m+1}$ of $C_m$, and
$4t=(1+t)^2-(1-t)^2$ give
\[
 R_{r,n}(y)=4^r d!\,i^{-d}E_{k,n}((iy-1)/2)
       =\sum_{0\le h\le\min(r,\lfloor d/2\rfloor)}
           (d)_{2h}\left(\sum_{j=h}^r 4^{r-j}g_j\binom jh\right)Q_{d-2h}(y).
\]
Here $(d)_a=d(d-1)\cdots(d-a+1)$. The leading coefficient is positive.

For $n\ge r+2$, apply the three-term recurrence in both directions to
write
\[
 R_{r,n}=A Q_n-B Q_{n-1},\qquad R_{r+1,n}=C Q_n-D Q_{n-1}.
\]
This expresses both polynomials using the pair $Q_n,Q_{n-1}$, as in the
repository's graph-tower section. In the first pair the degrees in $y$ are
at most $r,r-1$, respectively, and the respective parities are those of
$r,r-1$. For the second pair replace $r$ by $r+1$.
Define the two quantities
\begin{align*}
 H&=\text{coefficient matrix of }
      \frac{A(u)B(v)-B(u)A(v)}{u-v}
      \quad\text{in }1,u,\ldots,u^{r-1}\text{ and }1,v,\ldots,v^{r-1},\\
 T(y)&=A(y)D(y)-B(y)C(y).
\end{align*}
For $r>0$ the tests are $H$ positive definite and $T(y)>0$ on the real
axis. When $r=0$ take $A=1,B=0$ and omit the matrix test.

We recall why these tests suffice. In the upper half-plane, evaluating
$H$ at $(u,\bar u)$ shows that $A$ never vanishes there and
$\operatorname{Im}(B/A)<0$ for $r>0$. Meanwhile
$\operatorname{Im}(Q_n/Q_{n-1})>0$ by the recurrence, beginning with
$Q_1/Q_0=y$. Thus all roots of $R_{r,n}$ are real. On the real axis
$W=A'B-AB'>0$ by the diagonal of $H$ (use $W=0$ for $r=0$).
Also $W_n=Q_n'Q_{n-1}-Q_nQ_{n-1}'>0$, since
$W_1=1$ and $W_n=Q_{n-1}^2+(n-1)^2W_{n-1}$.
At any root $s$ of $R_{r,n}$, both pairs are nonzero and
$(Q_n,Q_{n-1})(s)=\lambda (B,A)(s)$ for some real $\lambda\ne0$.
Consequently
\[
 R_{r,n}'(s)=\frac{\lambda^2 W(s)+W_n(s)}{\lambda},\qquad
 \frac{R_{r+1,n}(s)}{R_{r,n}'(s)}
   =-\frac{\lambda^2 T(s)}{\lambda^2 W(s)+W_n(s)}<0.
\]
Each gap between consecutive roots therefore contains a root of
$R_{r+1,n}$. Its positive leading coefficient and extra degree give a
root in each exterior gap as well. This accounts for every root and proves
simplicity.

\emph{Certification over the entire half-line of integers.}
The bundled program \texttt{split\_interlacing.py} carries out the
following rational and integer calculations for $r=0,\ldots,19$.
For coordinates of width $s=r$ or $r+1$, put
$S_s(n)=\prod_{i=1}^{s-1}(n-i)^2$. Multiplication by $S_s$ clears the
backward recurrence denominators. To see the degree bound, note that the
forward coordinates have degree at most $2s$ in $n$. At relative index
$-t$, the backward coordinates have denominator
$\prod_{i=1}^{t-1}(n-i)^2$ and numerator degree at most $2t$, by
induction in $t$. Thus either coordinate after clearing has degree at
most $4s$. The weight contributes at most $2s+s$ more from the falling
factorial and $g_j$. Each cleared coordinate has degree at most $7s$.

At $7s+4$ consecutive integers starting at $b=r+2$, exact forward
differences certify that each coefficient of the first coordinate
agrees with a polynomial in $n$ of degree at most $s$, and each
coefficient of the second has degree at most $s+1$. Multiplying the
interpolant by $S_s$ makes this an identity by the preceding degree
bound (also for $s=0$). Each entry of $H$ therefore has degree at most
$2r+1$ in $n$, and $T(y)=V(n,y^2)$ has degree at most $2r+2$ in $n$.

Parity splits $H$ into an even and an odd block. For each leading minor
$f$ of size $a$ in either block, the program computes $f(b)>0$ and
$\Delta^j f(b)\ge0$ through degree $(2r+1)a$, by integer elimination.
Newton's identity $f(b+h)=\sum_j\binom hj\Delta^j f(b)$ proves $f(b+h)>0$
for every integer $h\ge0$, so Sylvester's criterion proves the matrix
test. For $0\le j\le2r+2$, the program forms the polynomial
$\Delta_n^jV(b,z)$ exactly. Each row is either zero or has a positive
constant term and is positive for $z\ge0$, as verified by its
coefficients or an exact Sturm root count. The zeroth row is positive.
Newton's identity proves $T(y)>0$ for every real $y$ and every integer
$n\ge b$.

For $1\le n<b$, the program computes both full rotated polynomials.
The negative remainders in their Euclidean algorithm strictly drop in
degree by one and have positive leading coefficients through the final
constant. Backwards induction, by evaluating at consecutive roots and at
both infinities, gives strict interlacing. All finite checks are recorded
in \texttt{interlacing\_checks.txt}. The full run took 48 seconds with Python
3.12.14 and SymPy 1.14.0.
\end{proof}
